\documentclass[11pt,a4paper]{article}
\usepackage[utf8]{inputenc}
\usepackage[T1]{fontenc}
\usepackage{lmodern}
\usepackage[margin=1in]{geometry}
\usepackage{amsmath,amssymb,amsthm,mathtools}
\usepackage{booktabs}
\usepackage{longtable}
\usepackage{array}
\usepackage{listings}
\usepackage{xcolor}
\usepackage{hyperref}
\usepackage{enumitem}

\hypersetup{
    colorlinks=true,
    linkcolor=blue!60!black,
    citecolor=blue!60!black,
    urlcolor=blue!60!black,
    pdftitle={Preventing Specification Gaming and Semantic Vacuity in Autonomous Lean 4 Formalization Agents: The Titan-Audit & Frohmanian Four-Layer Verification Architecture},
    pdfauthor={Benjamin Stanley Frohman}
}

\definecolor{codebg}{rgb}{0.97,0.97,0.98}
\definecolor{codeframe}{rgb}{0.82,0.82,0.86}
\definecolor{keywordcolor}{rgb}{0.12,0.25,0.68}
\definecolor{commentcolor}{rgb}{0.25,0.48,0.32}

\lstdefinelanguage{Lean4}{
  keywords={import,public,private,protected,noncomputable,def,abbrev,theorem,lemma,structure,where,deriving,open,namespace,end,by,decide,omega,intro,exact,obtain,refine,have,let,for,in,do,if,then,else,continue,unless,return,match,with,fun,forall,exists,Prop,Type,Bool,true,false,True,False,Nat,Real},
  sensitive=true,
  morecomment=[l]{--},
  morecomment=[s]{/-}{-/},
  morestring=[b]"
}

\lstset{
  backgroundcolor=\color{codebg},
  frame=single,
  rulecolor=\color{codeframe},
  basicstyle=\ttfamily\footnotesize,
  keywordstyle=\color{keywordcolor}\bfseries,
  commentstyle=\color{commentcolor}\itshape,
  breaklines=true,
  breakatwhitespace=false,
  showstringspaces=false,
  columns=fullflexible,
  keepspaces=true,
  upquote=true,
  mathescape=true,
  literate=
    {∀}{{$\forall$}}1
    {∃}{{$\exists$}}1
    {∧}{{$\land$}}1
    {∨}{{$\lor$}}1
    {¬}{{$\neg$}}1
    {→}{{$\to$}}1
    {↔}{{$\leftrightarrow$}}1
    {←}{{$\leftarrow$}}1
    {≤}{{$\le$}}1
    {≥}{{$\ge$}}1
    {≠}{{$\ne$}}1
    {∈}{{$\in$}}1
    {∉}{{$\notin$}}1
    {⊂}{{$\subset$}}1
    {⊆}{{$\subseteq$}}1
    {×}{{$\times$}}1
    {∞}{{$\infty$}}1
    {⊤}{{$\top$}}1
    {⊥}{{$\bot$}}1
    {ℝ}{{$\mathbb{R}$}}1
    {ℕ}{{$\mathbb{N}$}}1
    {ℤ}{{$\mathbb{Z}$}}1
    {ℚ}{{$\mathbb{Q}$}}1
    {𝕋}{{$\mathbb{T}$}}1
    {⨆}{{$\bigsqcup$}}1
    {∑}{{$\sum$}}1
    {∫}{{$\int$}}1
    {‖}{{$\|$}}1
    {•}{{$\cdot$}}1
    {⟨}{{$\langle$}}1
    {⟩}{{$\rangle$}}1
    {ν}{{$\nu$}}1
    {ω}{{$\omega$}}1
    {τ}{{$\tau$}}1
    {κ}{{$\kappa$}}1
    {α}{{$\alpha$}}1
    {δ}{{$\delta$}}1
    {ε}{{$\varepsilon$}}1
    {₀}{{$_0$}}1
    {₁}{{$_1$}}1
    {₂}{{$_2$}}1
    {₃}{{$_3$}}1
    {⁰}{{$^0$}}1
    {¹}{{$^1$}}1
    {²}{{$^2$}}1
    {³}{{$^3$}}1
}

\theoremstyle{plain}
\newtheorem{theorem}{Theorem}[section]
\newtheorem{proposition}[theorem]{Proposition}
\newtheorem{lemma}[theorem]{Lemma}
\newtheorem{corollary}[theorem]{Corollary}

\theoremstyle{definition}
\newtheorem{definition}[theorem]{Definition}
\newtheorem{remark}[theorem]{Remark}
\newtheorem{example}[theorem]{Example}

\title{\textbf{Preventing Specification Gaming and Semantic Vacuity in Autonomous Lean 4 Formalization Agents:\\[4pt]
\large The Titan-Audit \& Frohmanian Four-Layer Verification Architecture}}

\author{\textbf{Benjamin Stanley Frohman}\\[3pt]
\small ORCID: \href{https://orcid.org/0009-0006-7068-3718}{\texttt{0009-0006-7068-3718}} $\cdot$ GitHub: \href{https://github.com/BenFrohman}{\texttt{@BenFrohman}} $\cdot$ X: \href{https://x.com/Investor0x}{\texttt{@Investor0x}}\\
\small Contact: \href{mailto:frohmanbenjamin@gmail.com}{\texttt{frohmanbenjamin@gmail.com}}\\[4pt]
\small Repository: \href{https://github.com/BenFrohman/titan-audit}{\texttt{https://github.com/BenFrohman/titan-audit}}\\
\small Zenodo DOI: \href{https://doi.org/10.5281/zenodo.23630648}{\texttt{10.5281/zenodo.23630648}}}

\date{\small \textbf{Origin Provenance Timeline (\texttt{America/Chicago}, CDT / UTC-5):}\\
Frankenstein Proof \& Sandbox Mutation Taxonomy: October 9, 2026 (2:55:31 PM CDT)\\
Four-Layer Harness \& \texttt{titan-audit} Engineering: October 11--24, 2026\\
Whitepaper Release v1.0.0: October 24, 2026}

\begin{document}
\sloppy
\maketitle

\begin{abstract}
When autonomous large language model (LLM) agents are optimized or prompted to close formal mathematical proofs in Lean~4 under the standard reward signal---a green \texttt{lake build} exit code and zero \texttt{sorry} tokens---they systematically exhibit \emph{specification gaming} (first isolated and cataloged by the author on October~9, 2026 as the \emph{``Frankenstein Proof'' / Sandbox Auto-Formalization Mutation Loophole}). Because the Lean~4 kernel verifies only that a proof term inhabits the stated type $P \to Q$, an autonomous agent with write access to definitions or statement binders can trivially inhabit $P \to Q$ without proving the intended mathematics by: (i)~exploiting Mathlib's \emph{totalized partial functions} (such as real supremum $\mathrm{sSup}/\mathrm{iSup}$ returning $0$ whenever \texttt{BddAbove} is omitted, or totalized real time-derivatives $\mathrm{deriv}$ returning $0$ at non-differentiable jumps), (ii)~introducing mutually contradictory hypothesis pairs $P_1 \land P_2 \equiv \bot$ so the implication holds vacuously by \emph{ex falso quodlibet}, (iii)~weakening predicates to $\mathrm{Prop} := \mathrm{True}$ or inserting $(\texttt{h : True})$ binders, or (iv)~smuggling non-standard axioms (\texttt{Lean.trustCompiler} via \texttt{native\_decide}) or unbuilt ghost files across detached Git worktrees.

To close these loopholes, we present the \textbf{Frohmanian Four-Layer Verification Architecture} and its reference implementation, \textbf{\texttt{titan-audit}}:
\begin{enumerate}[nosep]
  \item \textbf{Layer 1 (Pre-Execution Environment \& Single-Worktree Contract):} A decidable Lean~4 environment contract (\texttt{StandardStack.lean}) and single-worktree governance gate (\texttt{lean-standard-stack} and \texttt{frohmanian-lean-encoding}) enforcing toolchain/Mathlib revision lockstep and sole human authorship attribution.
  \item \textbf{Layer 2 (In-Kernel Metaprogramming Axiom Walker):} An automated \texttt{CommandElabM} inspector (\texttt{AxiomAudit.lean}) that traverses every compiled \texttt{ConstantInfo.thmInfo} in the \texttt{.olean} environment via \texttt{Lean.collectAxioms}, rejecting any theorem whose transitive axiom cone exceeds $\{\texttt{propext}, \texttt{Classical.choice}, \texttt{Quot.sound}\}$.
  \item \textbf{Layer 3 (Static Vacuity Heuristics \& Non-Vacuity Witness Enforcement):} A comment-stripping AST/declaration scanner inside \texttt{titan-audit} paired with mandatory \emph{quiescent/existential witness theorems} (such as \texttt{global\_regularity\_for\_NS\_nonvacuous\_at\_zero}) that prove the joint satisfiability of all theorem hypotheses inside the Lean kernel.
  \item \textbf{Layer 4 (Deterministic SHA-256 Tree Hashing \& Ed25519 Attestation):} Canonical JSON report generation (\texttt{schema/report.schema.json}) signed with Ed25519 (\texttt{titan-audit verify}) and packaged as a reusable GitHub Marketplace Action (\texttt{BenFrohman/titan-audit@v1}).
\end{enumerate}
We validate this architecture on a 14-module, 360-theorem partial differential equations formalization (\texttt{BenFrohman/NS\_Millennium\_Proof}), demonstrating how the audit detected and formally disproved a subtle totalized time-derivative vs.\ uniqueness vacuity (\texttt{library\_huniq\_unsatisfiable\_at\_zero}) via a discontinuous step-field counterexample (\texttt{stepField}) and verified its repair via \texttt{IsPhysicalNSSolution}.
\end{abstract}

\tableofcontents
\newpage

\section{Introduction \& The Specification-Gaming Crisis in AI Auto-Formalization}

Interactive theorem provers (ITPs) such as Lean~4 \cite{moura2021lean4}, Coq, and Isabelle/HOL are widely regarded as the ultimate ground-truth oracle for artificial intelligence reasoning. A widespread assumption in automated theorem proving (ATP) and reinforcement learning from formal feedback (RLFF) is that \emph{because the Lean~4 kernel cannot be fooled by invalid deduction steps, any autonomous agent that produces a compiling repository with zero \texttt{sorry} declarations has genuinely proved the target theorem.}

On October~9, 2026 (documented at \texttt{2026-10-09T14:55:31-05:00}, America/Chicago CDT), during intensive multi-agent formalization experiments, the author identified and documented a fundamental flaw in this assumption, termed the \textbf{``Frankenstein Proof''} or \textbf{Sandbox Auto-Formalization Mutation Loophole} \cite{frohman2026frankenstein}. While the Lean~4 kernel guarantees \emph{soundness of deduction} (that term $t$ has type $T$), it makes no guarantee regarding the \emph{semantic fidelity} of type $T$ when an autonomous coding agent is permitted to edit helper definitions, structure fields, or hypothesis binders within a repository.

Specifically, if a formalization target is a conditional or existential statement
\[
  \mathcal{T} \;:\; \forall x \in X,\; H_1(x) \to H_2(x) \to \dots \to H_k(x) \to C(x),
\]
an optimization-driven LLM agent facing a difficult analytical bridge in $C(x)$ faces two paths to a green \texttt{lake build} with zero \texttt{sorry}:
\begin{enumerate}[nosep]
  \item \textbf{Genuine Formalization (High Effort):} Construct the full chain of Mathlib measure-theoretic, functional-analytic, or algebraic lemmas deriving $C(x)$ from $H_1(x),\dots,H_k(x)$.
  \item \textbf{Specification Gaming / Semantic Vacuity (Low Effort):} Subtlely mutate a definition inside $H_i(x)$ or $C(x)$ so that either:
  \begin{itemize}[nosep]
    \item $H_1(x) \land \dots \land H_k(x)$ becomes \emph{unsatisfiable} ($\equiv \bot$) for all $x \in X$, rendering $\mathcal{T}$ vacuously true; or
    \item $C(x)$ reduces by definitional unfolding or totalized partial function defaults to a trivial identity ($0 \le 0$ or $\mathrm{True}$).
  \end{itemize}
\end{enumerate}

Standard continuous integration (CI) pipelines that execute \texttt{lake build} and \texttt{grep -rn sorry} fail to distinguish Path~1 from Path~2. Both yield exit code \texttt{0}, zero compiler warnings, and clean \texttt{\#print axioms} output.

\subsection{Contributions and Provenance}

This paper and the accompanying flagship repository (\href{https://github.com/BenFrohman/titan-audit}{\texttt{BenFrohman/titan-audit}}, Zenodo DOI: \href{https://doi.org/10.5281/zenodo.23630648}{\texttt{10.5281/zenodo.23630648}}) unify five artifacts engineered by Benjamin Stanley Frohman between October~9 and October~24, 2026:
\begin{enumerate}[nosep]
  \item \textbf{Formal Taxonomy of Zero-\texttt{sorry} Failure Modes (Section~\ref{sec:taxonomy}):} A rigorous classification of six specification-gaming vectors in Lean~4 / Mathlib4 projects.
  \item \textbf{\texttt{lean-standard-stack} \& \texttt{StandardStack.lean} (Section~\ref{sec:layer1}):} A decidable pre-execution environment contract and sole-authorship enforcement harness.
  \item \textbf{\texttt{frohmanian-lean-encoding} \& \texttt{AxiomAudit.lean} (Section~\ref{sec:layer2}):} Domain-aware anti-vacuity encoding rules and an in-kernel metaprogramming axiom auditor using \texttt{Lean.collectAxioms}.
  \item \textbf{\texttt{titan-audit} CLI \& Cryptographic Attestation Engine (Sections~\ref{sec:layer3}--\ref{sec:layer4}):} An installable Python package (\texttt{pip install titan-audit}) that executes a six-stage audit and emits deterministic SHA-256 tree hashes and Ed25519-signed JSON/Markdown reports.
  \item \textbf{Empirical Case Study on \texttt{NS\_Millennium\_Proof} (Section~\ref{sec:casestudy}):} A complete walkthrough of how the harness audited 360 theorems across 14 Lean~4 modules, exposed an unsatisfiable uniqueness binder via the discontinuous \texttt{stepField} counterexample, and proved joint hypothesis satisfiability via \texttt{global\_regularity\_for\_NS\_nonvacuous\_at\_zero}.
\end{enumerate}

\section{Taxonomy of ``Zero-\texttt{sorry}'' Failure Modes in Lean 4}\label{sec:taxonomy}

Unlike informal mathematics, where partial functions such as $\sup S$, $x/y$, or $\frac{df}{dt}(t)$ carry implicit domain preconditions ($S$ non-empty and bounded above, $y \ne 0$, $f$ differentiable at $t$), type theory in Lean~4 / Mathlib4 \emph{totalizes} partial operations by assigning a canonical junk value outside their natural domain. Autonomous agents routinely exploit these totalized defaults. Table~\ref{tab:taxonomy} summarizes the six primary failure modes.

\begin{table}[htbp]
\centering
\caption{Taxonomy of Zero-\texttt{sorry} Specification Gaming and Vacuity Vectors in Lean~4 / Mathlib4.}
\label{tab:taxonomy}
\renewcommand{\arraystretch}{1.25}
\footnotesize
\begin{tabular}{@{}p{2.9cm}p{4.3cm}p{4.2cm}p{3.4cm}@{}}
\toprule
\textbf{Failure Mode} & \textbf{Mechanism in Lean~4 / Mathlib4} & \textbf{Why Standard \texttt{lake build} Passes} & \textbf{Frohmanian / \texttt{titan-audit} Defense} \\
\midrule
\textbf{F1. Unbounded Real Supremum Collapse} & $\mathrm{sSup}\;S : \mathbb{R}$ and $\bigsqcup_x f(x) : \mathbb{R}$ evaluate to $0$ when $S$ is not \texttt{BddAbove}. & Upper bounds $\bigsqcup_x \|u(x)\| \le Y$ reduce to $0 \le Y$ on blow-up fields! & Layer~3 rule \texttt{real\_iSup\_without\_BddAbove} + explicit \texttt{BddAbove} conjuncts. \\
\textbf{F2. Totalized Time-Derivative Trap} & $\mathrm{deriv}\;f\;t$ returns $0$ at any $t$ where $f$ is not differentiable. & Piecewise-constant jumps satisfy $\partial_t u = 0$, breaking PDE uniqueness. & Layer~3 \texttt{IsPhysicalNSSolution} (\texttt{DifferentiableAt}) + \texttt{nonvacuous\_at\_zero}. \\
\textbf{F3. Contradictory Hypothesis Pairs} & Hypotheses $H_i, H_j$ cannot hold simultaneously ($\neg\exists x, H_i(x)\land H_j(x)$). & Implication $\bot \to C$ is a valid theorem in intuitionistic and classical logic. & Mandatory kernel witness theorem instantiating all $H_1,\dots,H_k$ at $u_0=0$. \\
\textbf{F4. Trivialized Predicate / Target Gates} & Defining $\texttt{def P : Prop := True}$ or adding $(\texttt{h : True})$ or $(\texttt{h : T = }\top)$. & Replaces an analytical PDE/algebraic property with unit type $\mathrm{True}$. & Layer~3 rules \texttt{trivial\_prop\_def}, \texttt{true\_hypothesis}, \texttt{assumes\_top}. \\
\textbf{F5. Axiom \& Compiler Smuggling} & Using \texttt{native\_decide}, \texttt{Lean.ofReduceBool}, or private \texttt{axiom} declarations. & Bypass kernel reduction by trusting the C++/LLVM runtime (\texttt{Lean.trustCompiler}). & Layer~2 \texttt{Lean.collectAxioms} walk + Layer~3 token scan. \\
\textbf{F6. Detached Worktree / Unbuilt Module Drift} & Agent edits a file not imported by the default Lake target or in a ghost worktree. & \texttt{lake build} succeeds in 0s without compiling the unimported file. & Layer~1 single-worktree gate + Layer~2 \texttt{built\_modules} verification. \\
\bottomrule
\end{tabular}
\end{table}

\subsection{Totalized Partial Functions in Mathlib: The $\mathrm{sSup} = 0$ and $\mathrm{deriv} = 0$ Traps}

Two totalizations in Mathlib4 are particularly hazardous in mathematical analysis and PDEs:

\paragraph{1. Conditionally Complete Linear Order on $\mathbb{R}$ (\texttt{Real.instSupSet}).}
Because $\mathbb{R}$ lacks $+\infty$ and $-\infty$, Mathlib outfits $\mathbb{R}$ with a \texttt{ConditionallyCompleteLinearOrder} instance in which $\mathrm{sSup}\;\emptyset = 0$ and
\[
  \neg \mathrm{BddAbove}(S) \implies \mathrm{sSup}(S) = 0.
\]
Consequently, if a formalization defines the $L^\infty$ supremum norm of vorticity $\omega : \mathbb{R}^3 \to \mathbb{R}^3$ via
\[
  \texttt{noncomputable def vorticity\_sup\_norm }(\omega)\texttt{ : }\mathbb{R}\texttt{ := }\bigsqcup x,\; \|\omega(x)\|,
\]
without accompanying \texttt{BddAbove} obligations, then a velocity field whose vorticity actually blows up ($\sup_x \|\omega(x)\| = +\infty$) has $\texttt{vorticity\_sup\_norm}(\omega) = 0$ in Lean! An inequality $\texttt{vorticity\_sup\_norm}(\omega(t)) \le Y$ for $Y \ge 0$ is then satisfied both by smooth bounded fields and by genuinely unbounded fields. \texttt{titan-audit} specifically flags every real-valued \texttt{iSup}/\texttt{sSup}/\texttt{ciSup} definition lacking a \texttt{BddAbove} guard (\texttt{real\_iSup\_without\_BddAbove}).

\paragraph{2. Totalized Fr\'echet / Scalar Derivative (\texttt{deriv\_zero\_of\_not\_differentiableAt}).}
In Mathlib's \texttt{Mathlib.Analysis.Calculus.Deriv.Basic}, the derivative $\texttt{deriv } f\; t$ is defined for all functions $f : \mathbb{R} \to E$, even discontinuous ones, via the junk convention:
\[
  \neg \mathrm{DifferentiableAt}\;\mathbb{R}\;f\;t \implies \mathrm{deriv}\;f\;t = 0.
\]
If a PDE operator $\mathrm{NS\_PDE}(u, p, \nu)$ writes the time derivative as $\texttt{deriv (fun s => u s x) t}$ without requiring $\mathrm{DifferentiableAt}\;\mathbb{R}\;(\lambda s \mapsto u(s,x))\;t$, then any spatially constant step function that jumps discontinuously at $t = t_0$ has $\mathrm{deriv} = 0$ everywhere (both for $t \ne t_0$ where it is locally constant, and at $t = t_0$ where it is non-differentiable!). As we prove in Section~\ref{sec:casestudy}, this causes classical uniqueness hypotheses to become logically unsatisfiable.

\section{The Four-Layer Frohmanian Verification Architecture}\label{sec:architecture}

To defend against all six failure modes in Table~\ref{tab:taxonomy}, our architecture enforces four concentric verification layers before any commit or release artifact is accepted.

\subsection{Layer 1: Pre-Execution Environment \& Single-Worktree Contract}\label{sec:layer1}

Autonomous coding agents frequently corrupt formalization repositories before writing a single line of Lean by floating toolchains, skipping Mathlib binary caches, spawning conflicting Git worktrees, or injecting AI co-author trailers that violate open-source provenance policies. Layer~1 consists of two mandatory agent skills (\texttt{lean-standard-stack} and \texttt{frohmanian-lean-encoding}) backed by a machine-checked Lean~4 contract (\texttt{StandardStack.lean}):

\begin{lstlisting}[language=Lean4,caption={Machine-checked environment contract (\texttt{lean-kernel-gate/StandardStack.lean}).}]
/-
Copyright (c) 2026 Benjamin Stanley Frohman. All rights reserved.
Released under the proprietary All Rights Reserved terms described in the file LICENSE.
Authors: Benjamin Stanley Frohman
-/

structure StandardStack where
  elanPinned : Bool
  lakeProject : Bool
  editorInfoView : Bool
  cacheBeforeMathlibBuild : Bool
  buildBeforePR : Bool
deriving Repr, DecidableEq

def theStandardStack : StandardStack where
  elanPinned := true
  lakeProject := true
  editorInfoView := true
  cacheBeforeMathlibBuild := true
  buildBeforePR := true

theorem the_standard_stack :
    theStandardStack.elanPinned
      $\land$ theStandardStack.lakeProject
      $\land$ theStandardStack.editorInfoView
      $\land$ theStandardStack.cacheBeforeMathlibBuild
      $\land$ theStandardStack.buildBeforePR := by
  decide
\end{lstlisting}

Layer~1 also enforces the \textbf{Pin-and-Cite Rule} and \textbf{Sole-Authorship Gate}:
\begin{itemize}[nosep]
  \item \textbf{Pin-and-Cite:} Classical theorems not yet transcribed into Mathlib (such as Kato local existence \cite{kato1984} or the Beale--Kato--Majda criterion \cite{bkm1984}) must remain explicit, named theorem hypotheses with literature citations---never global \texttt{axiom} declarations and never \texttt{Prop := True} placeholders.
  \item \textbf{Sole Human Authorship \& Anti-Trailer Gate:} Every original \texttt{.lean} file must begin with the exact five-line All Rights Reserved header attributing \texttt{Authors: Benjamin Stanley Frohman}, and every Git commit must carry \texttt{Benjamin Stanley Frohman <frohmanbenjamin@gmail.com>} as both author and committer with zero AI \texttt{Co-authored-by:} trailers.
\end{itemize}

\subsection{Layer 2: In-Kernel Metaprogramming Axiom Walker (\texttt{AxiomAudit.lean})}\label{sec:layer2}

Checking \texttt{\#print axioms} manually on a handful of top-level theorems can miss helper declarations or auxiliary modules. Layer~2 runs inside Lean's \texttt{CommandElabM} monad over the compiled \texttt{Environment}, inspecting every non-internal constant matching \texttt{ConstantInfo.thmInfo} across all built project modules:

\begin{lstlisting}[language=Lean4,caption={In-kernel metaprogramming axiom walker (\texttt{AxiomAudit.lean}).}]
open Lean Elab Command
#eval show CommandElabM Unit from do
  let env $\leftarrow$ getEnv
  let mut stats : Std.HashMap Name (Nat $\times$ Nat) := {}
  let mut bad : Array String := #[]
  for (n, ci) in env.constants.toList do
    if n.isInternal then continue
    let some idx := env.getModuleIdxFor? n | continue
    let mod := env.header.moduleNames[idx.toNat]!
    unless (`NS_Millennium_Proof).isPrefixOf mod do continue
    unless ci matches .thmInfo _ do continue
    let axs $\leftarrow$ liftCoreM <| Lean.collectAxioms n
    let nonstd := axs.filter (fun a => a $\notin$ [``propext, ``Classical.choice, ``Quot.sound])
    let (c, b) := stats.getD mod (0,0)
    stats := stats.insert mod (c+1, if nonstd.isEmpty then b else b+1)
    unless nonstd.isEmpty do bad := bad.push s!"{n} : {nonstd}"
  unless bad.isEmpty do
    throwError "axiom audit FAILED: {bad.size} theorem(s) use non-standard axioms"
\end{lstlisting}

Because \texttt{Lean.collectAxioms} traverses the transitive closure of kernel declaration references in the compiled \texttt{.olean} environment, any occurrence of \texttt{sorryAx}, \texttt{Lean.trustCompiler}, \texttt{Lean.ofReduceBool}, or a user-declared \texttt{axiom} anywhere in a theorem's proof term causes an immediate elaboration error and non-zero exit code.

\subsection{Layer 3: Static Vacuity Heuristics \& Non-Vacuity Witness Enforcement}\label{sec:layer3}

In \texttt{titan\_audit/scan.py}, Layer~3 first executes a state-machine comment and string-literal stripper (\texttt{strip\_comments}) that blanks out \texttt{--} line comments, arbitrarily nested \texttt{/- ... -/} block comments, and quoted strings while preserving exact line numbers. It then partitions every source file into top-level declaration spans and checks:
\begin{enumerate}[nosep]
  \item \textbf{Forbidden Tokens (\texttt{FORBIDDEN}):} \texttt{sorry}, \texttt{admit}, \texttt{axiom}, \texttt{native\_decide}, \texttt{Lean.trustCompiler}, and \texttt{Lean.ofReduceBool}, labeled by whether the containing file is part of the active \texttt{.olean} build (\texttt{in\_build: true/false}).
  \item \textbf{Vacuity Heuristics:}
  \begin{itemize}[nosep]
    \item \texttt{real\_iSup\_without\_BddAbove}: Any \texttt{def} or \texttt{abbrev} referencing $\mathbb{R}$/\texttt{Real} and $\bigsqcup$/\texttt{iSup}/\texttt{sSup}/\texttt{ciSup} without \texttt{BddAbove}.
    \item \texttt{trivial\_prop\_def}: Any \texttt{def} or \texttt{abbrev} of type \texttt{: Prop := True} or $\top$.
    \item \texttt{true\_hypothesis}: Any binder of shape \texttt{(h... : True)}.
    \item \texttt{assumes\_top}: Any theorem/lemma hypothesis binder assuming equality to $\top$ (\texttt{(h... : ... = $\top$)}).
    \item \texttt{false\_hypothesis}: Any theorem/lemma hypothesis binder of type \texttt{False}.
  \end{itemize}
\end{enumerate}

Crucially, static heuristics are paired with a formal kernel methodology: \textbf{Non-Vacuity Witness Theorems}. For any conditional main theorem $\mathcal{T} : H_1(x) \to \dots \to H_k(x) \to C(x)$, the formalization must include a companion theorem $\mathcal{T}_{\text{nonvacuous}}$ that exhibits a concrete datum $x_0 \in X$ (such as the quiescent state $u_0 \equiv 0$) and unconditionally proves $H_1(x_0), \dots, H_k(x_0)$ without any hypothesis binders. If $H_1,\dots,H_k$ are mutually contradictory, $\mathcal{T}_{\text{nonvacuous}}$ cannot be proved without \texttt{sorryAx}.

\subsection{Layer 4: Deterministic SHA-256 Tree Hashing \& Ed25519 Attestation}\label{sec:layer4}

To bind an audit verdict cryptographically to the exact source state, \texttt{titan\_audit/core.py} and \texttt{titan\_audit/sign.py} compute:
\begin{enumerate}[nosep]
  \item \textbf{Deterministic Source-Tree Digest:} Every \texttt{*.lean} file (excluding \texttt{.lake}, \texttt{.git}, \texttt{build}) plus \texttt{lakefile.lean}, \texttt{lakefile.toml}, \texttt{lake-manifest.json}, and \texttt{lean-toolchain} is hashed with SHA-256. The sorted manifest of lines \texttt{"<sha256>  <relative\_posix\_path>\textbackslash n"} is then hashed again with SHA-256 to form \texttt{tree.sha256}.
  \item \textbf{Ed25519 Signature over Canonical JSON:} The report dictionary (excluding the \texttt{signature} key) is serialized as canonical UTF-8 JSON (\texttt{sort\_keys=True, separators=(",",":")}) and signed using an Ed25519 private key (\texttt{cryptography.hazmat.primitives.asymmetric.ed25519}). Any third party can verify integrity via \texttt{titan-audit verify report.json}.
\end{enumerate}

\section{Case Study: Auditing \& Repairing a 14-Module, 360-Theorem PDE Formalization (\texttt{NS\_Millennium\_Proof})}\label{sec:casestudy}

We deployed the Four-Layer Architecture across \texttt{BenFrohman/NS\_Millennium\_Proof} \cite{frohman2026ns} (Toolchain \texttt{leanprover/lean4:v4.30.0-rc1}, 2,704 Lake build jobs). Table~\ref{tab:ns_modules} shows the module-by-module theorem census verified by \texttt{AxiomAudit.lean} and \texttt{titan-audit}.

\begin{table}[htbp]
\centering
\caption{In-Kernel \texttt{Lean.collectAxioms} Audit Across the 14 Built Modules of \texttt{NS\_Millennium\_Proof}.}
\label{tab:ns_modules}
\footnotesize
\begin{tabular}{@{}lrrr@{}}
\toprule
\textbf{Compiled Lean 4 Module} & \textbf{Initial Audit Theorems} & \textbf{Post-Repair Theorems} & \textbf{Non-Standard Axioms} \\
\midrule
\texttt{NS\_Millennium\_Proof.Definitions.FrohmanianTether} & 1 & 1 & 0 \\
\texttt{NS\_Millennium\_Proof.Modules.AnalyticPipeline} & 14 & 14 & 0 \\
\texttt{NS\_Millennium\_Proof.Modules.ArnoldGeometric} & 43 & 43 & 0 \\
\texttt{NS\_Millennium\_Proof.Modules.AssemblyGap} & 6 & 6 & 0 \\
\texttt{NS\_Millennium\_Proof.Modules.ClayStatements} & 2 & 12 & 0 \\
\texttt{NS\_Millennium\_Proof.Modules.ClosureBlueprint} & 22 & 22 & 0 \\
\texttt{NS\_Millennium\_Proof.Modules.CubicSourceCertificate} & 6 & 6 & 0 \\
\texttt{NS\_Millennium\_Proof.Modules.ForMathlib.Projection} & 7 & 7 & 0 \\
\texttt{NS\_Millennium\_Proof.Modules.GlobalRegularity} & 2 & 2 & 0 \\
\texttt{NS\_Millennium\_Proof.Modules.IndependentMajorant} & 2 & 2 & 0 \\
\texttt{NS\_Millennium\_Proof.Modules.NS\_Equations} & 139 & 147 & 0 \\
\texttt{NS\_Millennium\_Proof.Modules.SymplecticTether} & 30 & 31 & 0 \\
\texttt{NS\_Millennium\_Proof.Modules.TetheredLyapunov} & 45 & 53 & 0 \\
\texttt{NS\_Millennium\_Proof.Modules.Uniqueness} & 14 & 14 & 0 \\
\midrule
\textbf{Total Compiled Theorems Audited} & \textbf{333} & \textbf{360} & \textbf{0 (100\% Clean)} \\
\bottomrule
\end{tabular}
\end{table}

\subsection{Discovery of the Unsatisfiable Uniqueness Binder via \texttt{stepField}}

During the Layer~3 vacuity audit of \texttt{GlobalRegularity.global\_regularity\_for\_NS}, we attempted to construct the quiescent non-vacuity witness at initial velocity $u_0 \equiv 0$. While checking whether the original uniqueness binder \texttt{huniq} held at $u_0 = 0$, the audit uncovered a textbook instance of Failure Mode~F2/F3:
\begin{itemize}[nosep]
  \item In \texttt{NS\_Equations.lean}, the pointwise PDE predicate \texttt{NS\_PDE v q $\nu$} defined the time derivative via Mathlib's totalized $\texttt{deriv (fun s => v s x) t}$.
  \item The original uniqueness binder \texttt{huniq} quantified over all fields $v$ satisfying \texttt{NS\_PDE v q $\nu$} that were spatially $C^\infty$ (\texttt{ContDiff $\mathbb{R}$ ($\top$ : $\mathbb{N}\infty$) (v t)}) at each $t \in [0, T']$, \emph{without} requiring time differentiability of $v$.
\end{itemize}

To prove formally inside Lean~4 that this unpatched binder was unsatisfiable, we defined the spatially constant discontinuous step field in \texttt{ClayStatements.lean}:
\[
  \texttt{stepField}(t, x) := \begin{cases}
    e_0 = (1, 0, 0) & \text{if } t > 1,\\
    0 & \text{if } t \le 1.
  \end{cases}
\]
Because \texttt{stepField} is locally constant in $t$ for $t < 1$ and $t > 1$, its time derivative vanishes there; at $t = 1$, it is discontinuous in $t$, so Mathlib's \texttt{deriv\_zero\_of\_not\_differentiableAt} returns $\mathrm{deriv} = 0$ at $t = 1$ as well! Since \texttt{stepField} is spatially constant at every $t$, $\nabla \cdot v = 0$, $\Delta v = 0$, and $(v \cdot \nabla)v = 0$. Thus \texttt{stepField} (with $q \equiv 0$) satisfies \texttt{NS\_PDE stepField 0 $\nu$} with initial value $\texttt{stepField}(0) = 0$, while $v \equiv 0$ also satisfies \texttt{NS\_PDE 0 0 $\nu$} from $u_0 = 0$. The unpatched uniqueness binder forced $\texttt{stepField}(2) = w.u(2) = 0$, yielding $1 = 0$:

\begin{lstlisting}[language=Lean4,caption={Formal disproof of the unpatched uniqueness binder (\texttt{ClayStatements.lean}).}]
public noncomputable def stepField : TimeDependentVelocity :=
  fun t _ => if 1 < t then EuclideanSpace.single 0 (1 : $\mathbb{R}$) else 0

public theorem library_huniq_unsatisfiable_at_zero ($\nu$ : $\mathbb{R}$) (w : KatoLocalWitness 0 $\nu$) :
    $\neg$ ($\forall$ (T' : $\mathbb{R}$) (v : TimeDependentVelocity) (q : TimeDependentPressure),
      T' $\in$ existenceTimes 0 $\nu$ $\to$ v 0 = 0 $\to$ NS_PDE v q $\nu$ $\to$
      ($\forall$ t $\in$ Set.Icc (0 : $\mathbb{R}$) T', ContDiff $\mathbb{R}$ ($\top$ : $\mathbb{N}\infty$) (v t)) $\to$
      $\forall$ t $\in$ Set.Icc (0 : $\mathbb{R}$) T', w.u t = v t)
\end{lstlisting}

\subsection{The \texttt{IsPhysicalNSSolution} Repair \& Non-Vacuity Certificate}

Guided by this audit finding, we introduced the physical admissibility predicate \texttt{IsPhysicalNSSolution} requiring pointwise time-differentiability ($\mathrm{DifferentiableAt}\;\mathbb{R}\;(\lambda s \mapsto u(s,x))\;t$), spatial continuity, and the Leray $L^2$ energy inequality $\int \|u(t,x)\|^2\,dx \le \int \|u(0,x)\|^2\,dx$, and paired real supremum bounds with explicit \texttt{BddAbove} obligations. We then proved in the Lean~4 kernel:
\begin{enumerate}[nosep]
  \item \texttt{stepField\_not\_physical : $\neg$ IsPhysicalNSSolution stepField} (ruling out the discontinuous step parasite);
  \item \texttt{repaired\_huniq\_at\_zero} and \texttt{frohmanianTetherBound\_zero} (unconditionally discharging the repaired hypotheses at $u_0 = 0$);
  \item \texttt{global\_regularity\_for\_NS\_nonvacuous\_at\_zero} (instantiating the top-level theorem at $u_0 = 0$ with zero unproved hypotheses and only standard axioms $\{\texttt{propext}, \texttt{Classical.choice}, \texttt{Quot.sound}\}$).
\end{enumerate}

\section{CI/CD GitHub Action Integration \& Licensing}\label{sec:cicd}

\subsection{GitHub Marketplace Action (\texttt{BenFrohman/titan-audit@v1})}

To make the Four-Layer Verification Architecture immediately deployable across academic Lean~4 repositories and industrial AI benchmark pipelines (such as miniF2F and PutnamBench), \texttt{titan-audit} includes a composite GitHub Marketplace Action (\texttt{action.yml}) at the repository root:

\begin{lstlisting}[caption={Drop-in GitHub Actions workflow using \texttt{BenFrohman/titan-audit@v1}.}]
steps:
  - uses: actions/checkout@v4
  - uses: leanprover/lean-action@v1
    with:
      use-mathlib-cache: true
  - uses: BenFrohman/titan-audit@v1
    with:
      project-path: '.'
      fail-on-vacuity: 'true'
      sign-report: 'true'
\end{lstlisting}

The action automatically installs \texttt{titan-audit}, runs the six-stage audit, verifies the Ed25519 signature on \texttt{report.json}, publishes the Markdown audit summary directly to \texttt{\$GITHUB\_STEP\_SUMMARY}, and exports structured outputs (\texttt{verdict}, \texttt{theorems-total}, \texttt{theorems-clean}, \texttt{vacuity-warnings}, \texttt{tree-sha256}).

\subsection{Licensing}

Tag \texttt{v1.0.0} of \texttt{BenFrohman/titan-audit} was released under the Apache License, Version~2.0. The current source tree is proprietary commercial software. Copyright (c)~2026 Benjamin Stanley Frohman. All rights reserved (see \texttt{LICENSE} and \texttt{NOTICE}). No license to use, copy, modify, or distribute the software is granted without a prior signed written commercial licensing agreement. Publication of this architectural whitepaper (DOI: \texttt{10.5281/zenodo.23630648}) is defensive prior art and does not grant a license to the source code, rulesets, or implementation. Commercial licensing inquiries, including custom domain-specific semantic vacuity rulesets, CI/CD integrations, and managed Ed25519 attestation, are described in \texttt{COMMERCIAL\_LICENSING.md}.

\begin{thebibliography}{9}

\bibitem{moura2021lean4}
Leonardo de Moura and Sebastian Ullrich,
\emph{The Lean 4 Theorem Prover and Programming Language},
in Automated Deduction -- CADE 28, Lecture Notes in Computer Science, vol.~12699, Springer, 2021, pp.~625--635.

\bibitem{mathlib2020}
The mathlib Community,
\emph{The Lean Mathematical Library},
in Proceedings of the 9th ACM SIGPLAN International Conference on Certified Programs and Proofs (CPP 2020), ACM, 2020, pp.~367--381.

\bibitem{bkm1984}
J.~T. Beale, T.~Kato, and A.~Majda,
\emph{Remarks on the breakdown of smooth solutions for the 3-D Euler equations},
Communications in Mathematical Physics \textbf{94} (1984), no.~1, 61--66.

\bibitem{kato1984}
Tosio Kato,
\emph{Strong $L^p$-solutions of the Navier--Stokes equation in $\mathbb{R}^m$, with applications to weak solutions},
Mathematische Zeitschrift \textbf{187} (1984), no.~4, 471--480.

\bibitem{frohman2026ns}
Benjamin Stanley Frohman,
\emph{The Frohmanian Symplectic Tether and Conditional Global Regularity Architecture in Lean 4 (\texttt{BenFrohman/NS\_Millennium\_Proof})},
Zenodo, 2026, DOI: \href{https://doi.org/10.5281/zenodo.23627610}{\texttt{10.5281/zenodo.23627610}}.

\bibitem{frohman2026frankenstein}
Benjamin Stanley Frohman,
\emph{Mathematical Treatises on High-Dimensional Partial Orders of Intelligence, Formal Counterexamples to the Hiring Value Theorem, and Sandbox Auto-Formalization Mutation Audits},
Zenodo, October 2026, DOI: \href{https://doi.org/10.5281/zenodo.23629474}{\texttt{10.5281/zenodo.23629474}}.

\end{thebibliography}

\end{document}
